\clearpage
\section{교대군 A5의 켤레류와 단순성}
\label{sec:a5-conjugacy-simplicity}

\begin{learninggoal}
$A_5$의 원소를 순환형에 따라 분류하고 켤레류의 크기를 계산한다. 정상부분군은 켤레류들의 합집합이라는 사실을 이용하여 $A_5$가 단순군임을 증명하고, 이 구조가 오차방정식의 비가해성과 어떻게 연결되는지 이해한다.
\end{learninggoal}

오차방정식에서 처음 나타나는 본질적인 장애물은 단순히 군의 위수가 크다는 데 있지 않다. 핵심은 군 안에 더 이상 아벨 몫으로 분해할 수 없는 비가환 핵이 존재한다는 점이다. 그 최소 표준 예가 위수 $60$의 교대군 $A_5$이다.

\subsection{교대군 A5의 원소 유형}

$A_5$의 원소는 $5$개 문자의 짝순열이다. 서로소 순환분해를 사용하면 가능한 비항등 순환형은 다음 세 가지뿐이다.
\[
  3+1+1,\qquad 2+2+1,\qquad 5.
\]
즉 비항등원은 삼순환, 서로소인 두 전치의 곱, 또는 오순환이다.

\begin{proposition}[$A_5$의 원소 개수]
\label{prop:a5-element-counts}
$A_5$의 원소는 다음과 같이 분포한다.
\begin{center}
\renewcommand{\arraystretch}{1.2}
\begin{tabular}{lccc}
\toprule
순환형 & 대표원소 & 개수 & 원소의 위수 \\
\midrule
$1^5$ & $1$ & $1$ & $1$ \\
$3+1+1$ & $(1\,2\,3)$ & $20$ & $3$ \\
$2+2+1$ & $(1\,2)(3\,4)$ & $15$ & $2$ \\
$5$ & $(1\,2\,3\,4\,5)$ & $24$ & $5$ \\
\bottomrule
\end{tabular}
\end{center}
특히 $|A_5|=1+20+15+24=60$이다.
\end{proposition}

\begin{proof}
삼순환은 움직일 세 문자를 고르는 방법이 $\binom53$가지이고, 선택된 세 문자 위에는 서로 역인 두 삼순환이 있으므로
\[
  \binom53\cdot2=20
\]
개이다.

서로소인 두 전치의 곱은 먼저 고정할 한 문자를 $5$가지로 고르고, 나머지 네 문자를 두 쌍으로 나누는 방법이 $3$가지이므로 $15$개이다.

오순환의 개수는 잘 알려진 대로
\[
  (5-1)!=24
\]
개이다. 이 세 유형은 모두 짝순열이고, 그 개수의 합에 항등원을 더하면 $60=5!/2$가 된다.
\end{proof}

\subsection{켤레류의 크기}

정상부분군은 켤레에 닫혀 있으므로 켤레류의 구조를 알아야 한다. 군 $G$에서 원소 $g$의 켤레류 크기는
\[
  |\operatorname{Cl}_G(g)|=[G:C_G(g)]
\]
이다.

\begin{proposition}[$A_5$의 켤레류]
\label{prop:a5-conjugacy-classes}
$A_5$의 켤레류 크기는
\[
  1,\qquad 20,\qquad 15,\qquad 12,\qquad 12
\]
이다. 구체적으로 삼순환들은 하나의 크기 $20$인 켤레류를 이루고, 이중전치들은 하나의 크기 $15$인 켤레류를 이루며, 오순환 $24$개는 크기 $12$인 두 켤레류로 갈라진다.
\end{proposition}

\begin{proof}
$\alpha=(1\,2\,3)$라 하자. $S_5$에서 $\alpha$의 중심화군은
\[
  \langle\alpha\rangle\times\langle(4\,5)\rangle
\]
이고 위수는 $6$이다. 그 가운데 짝순열은 $\langle\alpha\rangle$의 세 원소뿐이므로
\[
  |C_{A_5}(\alpha)|=3,
  \qquad
  |\operatorname{Cl}_{A_5}(\alpha)|=60/3=20.
\]
따라서 모든 삼순환이 하나의 켤레류를 이룬다.

$\beta=(1\,2)(3\,4)$의 $S_5$-중심화군은 위수 $8$이고, 그 짝순열 부분은
\[
  \{1,(1\,2)(3\,4),(1\,3)(2\,4),(1\,4)(2\,3)\}
\]
인 클라인 네 군이다. 따라서
\[
  |C_{A_5}(\beta)|=4,
  \qquad
  |\operatorname{Cl}_{A_5}(\beta)|=60/4=15.
\]

마지막으로 $\gamma=(1\,2\,3\,4\,5)$의 $S_5$-중심화군은 $\langle\gamma\rangle$이고 위수는 $5$이다. 오순환의 모든 거듭제곱은 짝순열이므로
\[
  C_{A_5}(\gamma)=\langle\gamma\rangle,
  \qquad
  |\operatorname{Cl}_{A_5}(\gamma)|=60/5=12.
\]
오순환은 모두 $24$개이므로 $A_5$ 안에서는 두 개의 크기 $12$인 켤레류로 갈라진다.
\end{proof}

\begin{remark}
$S_5$에서는 같은 순환형의 원소가 항상 서로 켤레이다. 그러나 $A_5$에서는 오순환의 $S_5$-켤레류가 둘로 갈라진다. 이는 $A_n$의 켤레류가 $S_n$의 켤레류보다 더 미세할 수 있음을 보여 주는 첫 중요한 예이다.
\end{remark}

\subsection{교대군 A5의 단순성}

\begin{definition}[단순군]
비자명한 군 $G$가 정상부분군을
\[
  1\quad\text{과}\quad G
\]
밖에 갖지 않을 때 $G$를 \emph{단순군}(simple group)이라 한다.
\end{definition}

\begin{theorem}[$A_5$의 단순성]
\label{thm:a5-simple}
교대군 $A_5$는 비가환 단순군이다.
\end{theorem}

\begin{proofidea}
정상부분군은 항등원을 포함하는 켤레류들의 합집합이다. 따라서 그 위수는
\[
  1,15,20,12,12
\]
중 일부를 더한 값이어야 한다. 동시에 라그랑주 정리에 의해 $60$의 약수여야 한다. 가능한 합을 조사하면 $1$과 $60$밖에 남지 않는다.
\end{proofidea}

\begin{proof}
$N\triangleleft A_5$라 하자. 정상성 때문에 $N$은 $A_5$의 켤레류들의 합집합이며 반드시 항등원 켤레류를 포함한다. 따라서 $|N|$은
\[
  1+\text{선택된 일부의 }15,20,12,12
\]
꼴이다. 가능한 값은
\[
\begin{gathered}
1,13,16,21,25,28,33,36,40,45,48,60.
\end{gathered}
\]
한편 $|N|$은 $60$을 나누어야 한다. 위 목록에서 $60$의 약수인 것은 $1$과 $60$뿐이다. 따라서
\[
  N=1\qquad\text{또는}\qquad N=A_5.
\]
또한 $A_5$는 삼순환들이 서로 가환하지 않으므로 비가환이다.
\end{proof}

\begin{corollary}
\label{cor:a5-perfect-nonsolvable}
$A_5$는 완전군이며 가해군이 아니다. 즉
\[
  [A_5,A_5]=A_5.
\]
\end{corollary}

\begin{proof}
교환자부분군 $[A_5,A_5]$는 정상부분군이다. $A_5$가 비가환이므로 이 부분군은 $1$이 아니다. 정리 \ref{thm:a5-simple}에 의해 $[A_5,A_5]=A_5$이다. 따라서 도함열은 첫 단계 이후에도 줄어들지 않으므로 $A_5$는 가해군이 아니다.
\end{proof}

\begin{keyidea}
단순성은 ``부분군이 없다''는 뜻이 아니다. $A_5$에는 많은 부분군이 있지만, 전체 군 안에서 켤레에 안정적인 비자명한 중간층이 없다. 근호해법에 필요한 것은 바로 정상부분군의 연쇄이므로, 이 부재가 오차방정식의 핵심 장애물이 된다.
\end{keyidea}

\begin{checkpoint}
\begin{enumerate}[label=(\alph*)]
  \item $A_5$의 오순환 두 켤레류 가운데 $(1\,2\,3\,4\,5)$와 같은 류에 속하는 거듭제곱을 찾아라.
  \item 켤레류 크기 $1,20,15,12,12$의 합 가운데 $30$이 나타나지 않는 이유를 설명하라.
  \item 단순 아벨군은 정확히 어떤 군들인지 설명하고, $A_5$와 비교하라.
\end{enumerate}
\end{checkpoint}
